% ==============================================================================
% ================================= Aufgabe 1 =================================
% ==============================================================================

\chapter{Aufgabe 1: Datenbank}
\label{chap:aufgabe1}
\thispagestyle{fancy}

Im Rahmen der ersten Aufgabe wird die Datenbankstruktur für die Umfrageanwendung entworfen und implementiert. Die Datenbank besteht aus zwei Tabellen: einer Tabelle für die Nutzerverwaltung und einer Tabelle für die Speicherung der Umfrageantworten.

% ======================================= a =====================================

\section{Tabelle \texttt{user\_table} (\textit{1a})}

Die Tabelle \texttt{user\_table} dient der Verwaltung der Umfrageteilnehmer. Sie speichert die Namen der Nutzer und wird für die Zuordnung von Umfrageantworten zu bestimmten Nutzern benötigt.

\subsection{Tabellenstruktur}

Die Tabelle \texttt{user\_table} wurde mit folgender Struktur angelegt:

\begin{table}[h]
\centering
\begin{tabular}{|l|l|l|}
\hline
\textbf{Spalte} & \textbf{Datentyp} & \textbf{Besonderheit} \\
\hline
Id & INT & Primary Key, Auto Increment \\
Name & VARCHAR(255) & NOT NULL \\
\hline
\end{tabular}
\caption{Struktur der Tabelle \texttt{user\_table}}
\label{tab:user_table}
\end{table}

Die Spalte \texttt{Id} dient als eindeutiger Primärschlüssel und wird automatisch inkrementiert. Die Spalte \texttt{Name} speichert den Namen des Nutzers. Ein NULL-Wert ist hier nicht zugelassen, da jeder Nutzer einen Namen besitzen muss – bei anonymen Teilnahmen wird der Name \texttt{anonymous} verwendet.

\subsection{\acs{SQL}-Statement}

Das Anlegen der Tabelle erfolgte mit folgendem \acs{SQL}-Statement:

\lstinputlisting[
    language=sql, 
    caption={\acs{SQL}-Statement für \texttt{user\_table}}, 
    label={lst:sql_statement_user_table}
]{asset/code/Aufgabe1a.tex}

% ======================================= b =====================================

\section{Tabelle \texttt{survey\_table} (\textit{1b})}

Die Tabelle \texttt{survey\_table} speichert die eigentlichen Umfrageantworten. Jeder Datensatz repräsentiert eine vollständig ausgefüllte Umfrage eines Nutzers.

\subsection{Tabellenstruktur}

Die Tabelle \texttt{survey\_table} wurde mit folgender Struktur angelegt:

\begin{table}[h]
\centering
\begin{tabular}{|l|l|l|}
\hline
\textbf{Spalte} & \textbf{Datentyp} & \textbf{Besonderheit} \\
\hline
Id & INT & Primary Key, Auto Increment \\
UserId & INT & NOT NULL, Fremdschlüssel \\
Datum & DATE & NOT NULL \\
Q1 & VARCHAR(255) & NOT NULL \\
Q2 & VARCHAR(255) & NOT NULL \\
Q3 & VARCHAR(255) & NOT NULL \\
Q4 & VARCHAR(255) & NOT NULL \\
Q5 & VARCHAR(255) & NOT NULL \\
\hline
\end{tabular}
\caption{Struktur der Tabelle \texttt{survey\_table}}
\label{tab:survey_table}
\end{table}

Die Spalte \texttt{UserId} ist ein Fremdschlüssel, der auf den Primärschlüssel der Tabelle \texttt{user\_table} verweist. Dadurch wird sichergestellt, dass jede Umfrage einem existierenden Nutzer zugeordnet wird. Die Spalte \texttt{Datum} speichert das Datum der Umfrageteilnahme. Die Spalten \texttt{Q1} bis \texttt{Q5} halten die Antworten zu den fünf Umfragefragen.

Die Wahl des Datentyps \texttt{VARCHAR(\textit{255})} für die Antwortfelder ermöglicht sowohl kurze als auch längere Antwortformate und ist damit für verschiedene Fragetypen flexibel einsetzbar.

\subsection{\acs{SQL}-Statement}

Das Anlegen der Tabelle erfolgte mit folgendem \acs{SQL}-Statement:

\lstinputlisting[
    language=sql, 
    caption={\acs{SQL}-Statement für \texttt{survey\_table}}, 
    label={lst:sql_statement_survey_table}
]{asset/code/Aufgabe1b.tex}

% ======================================= c =====================================

\section{Befüllen der \texttt{user\_table} (\textit{1c})}

Gemäß der Aufgabenstellung wurde die Tabelle \texttt{user\_table} mit drei Datensätzen befüllt. Ein Datensatz repräsentiert den anonymen Teilnehmer, die beiden weiteren Datensätze repräsentieren namentlich registrierte Nutzer.

\subsection{INSERT-Statements}

Die folgenden \acs{SQL}-Statements wurden zur Befüllung der Tabelle ausgeführt:

\lstinputlisting[
    language=sql, 
    caption={INSERT-Statements für \texttt{user\_table}}, 
    label={lst:insert_statement_user_table}
]{asset/code/Aufgabe1c.tex}

\subsection{Ergebnis}

Nach der Ausführung enthält die Tabelle \texttt{user\_table} folgende Datensätze:

\begin{table}[h]
\centering
\begin{tabular}{|l|l|}
\hline
\textbf{Id} & \textbf{Name} \\
\hline
1 & anonymous \\
2 & Lara \\
3 & Luisa \\
\hline
\end{tabular}
\caption{Datensätze der Tabelle \texttt{user\_table}}
\label{tab:user_data}
\end{table}

\section{Export der Datenbank}

Nach dem Anlegen der Tabellen und dem Einfügen der Testdatensätze wurde die gesamte Datenbank über die Export-Funktion von phpMyAdmin als SQL-Datei gesichert. Die exportierte Datei enthält alle SQL-Statements aus den Aufgaben 1a, 1b und 1c:

\begin{itemize}
    \item \texttt{CREATE TABLE user\_table ...}
    \item \texttt{CREATE TABLE survey\_table ...}
    \item \texttt{INSERT INTO user\_table ...} (drei Datensätze)
    \item Fremdschlüssel-Constraints
    \item Tabellen-Indizes
\end{itemize}

Die Export-Datei wurde von phpMyAdmin als \texttt{umfrage\_db.sql} bereitgestellt und gemäß der Aufgabenstellung in \texttt{db.sql} umbenannt. Sie liegt nun im Projektverzeichnis unter folgendem Pfad:

\begin{forest}
    for tree={
        font=\ttfamily,
        grow'=0,
        child anchor=west,
        parent anchor=south,
        anchor=west,
        calign=first,
        edge path={
            \noexpand\path [draw, \forestoption{edge}] (!u.south west) +(10pt,0) |- node[fill,inner sep=1.5pt] {} (.child anchor)\forestoption{edge label};
        },
        before typesetting nodes={
            if n=1{
                insert before={[,phantom,calign with current]}
            }{}
        },
        fit=band,
        before computing xy={l=22pt},
    }
    [\textbf{C:\textbackslash xampp\textbackslash htdocs\textbackslash survey}
        [\textbf{inc\textbackslash}
            [\textbf{db.sql}]
        ]
    ]
\end{forest}

Die vollständige Datenbank kann jederzeit durch Einspielen der \texttt{db.sql}-Datei wiederhergestellt werden:

\section{Zusammenfassung}

Die Datenbank \texttt{umfrage\_db} besteht aus den zwei Tabellen \texttt{user\_table} und \texttt{survey\_table}, die über einen Fremdschlüssel miteinander verbunden sind. Die Tabelle \texttt{user\_table} wurde mit den erforderlichen Testdatensätzen befüllt.

Die physischen Datenbankdateien liegen im Verzeichnis \texttt{C:\textbackslash xampp\textbackslash mysql\textbackslash data\textbackslash umfrage\_db\textbackslash}.

Die Datenbank ist nun bereit für die Anbindung an die \acs{PHP}-Anwendung.

% ==============================================================================
% =================================== Aufgabe 1 ================================
% ==============================================================================